%%% preamble.tex -- shared preamble for Ben Webster's papers.
%%%
%%% Loaded from the driver with %%% preamble.tex -- shared preamble for Ben Webster's papers.
%%%
%%% Loaded from the driver with %%% preamble.tex -- shared preamble for Ben Webster's papers.
%%%
%%% Loaded from the driver with %%% preamble.tex -- shared preamble for Ben Webster's papers.
%%%
%%% Loaded from the driver with \input{preamble} immediately after
%%% \documentclass.  Project-specific macros belong in main.tex (or their own
%%% file), not here: this file is meant to be copied unchanged between papers.
%%%
%%%   1. Packages
%%%   2. Bibliography (biblatex + biber)
%%%   3. Cross-references (zref-clever + aliascnt)
%%%   4. Theorem environments
%%%   5. Equation numbering
%%%   6. Index of notation (anindex)
%%%   7. Standard macros
%%%   8. Editorial macros


%%%%%%%%%%%%%%%%%%%%%%%%%%%%%%%%%%%%%%%%%%%%%%%%%%%%%%%%%%%%%%%%%%%%%%%%%%%%%%
%%% 1. Packages
%%%%%%%%%%%%%%%%%%%%%%%%%%%%%%%%%%%%%%%%%%%%%%%%%%%%%%%%%%%%%%%%%%%%%%%%%%%%%%

% These two must be set before bm loads; without them bm + stmaryrd together
% exhaust the 16 available math families.
\newcommand{\hmmax}{0}
\newcommand{\bmmax}{0}

\usepackage[english]{babel}
\usepackage{csquotes}
\usepackage[letterpaper,top=2cm,bottom=2cm,left=3cm,right=3cm,marginparwidth=1.75cm]{geometry}

\usepackage{amsmath,amsfonts,amssymb,mathtools}
\usepackage{mathrsfs}                 % \mathscr
\usepackage{bm,stmaryrd,nicefrac}
\usepackage{xcolor,graphicx,xspace}
\usepackage{xparse,enumitem}
\usepackage{tikz,tikz-cd}
\usepackage{quiver}                   % local: quiver.sty
\usepackage{tabularx}

% Uncomment while drafting to print label names in the margin.
% \usepackage[notref,notcite]{showkeys}


%%%%%%%%%%%%%%%%%%%%%%%%%%%%%%%%%%%%%%%%%%%%%%%%%%%%%%%%%%%%%%%%%%%%%%%%%%%%%%
%%% 2. Bibliography
%%%%%%%%%%%%%%%%%%%%%%%%%%%%%%%%%%%%%%%%%%%%%%%%%%%%%%%%%%%%%%%%%%%%%%%%%%%%%%

\usepackage[maxbibnames=99,backend=biber,url=false,style=alphabetic,doi=true,eprint=false]{biblatex}
\renewbibmacro*{url+urldate}{%
  % nothing here
}
\renewbibmacro{in:}{}
\AtEveryBibitem{
  \clearfield{isbn}
  \clearlist{location}
  \clearfield{month}
  \clearfield{day}
  \clearfield{series}
  \clearfield{pages}
  \clearlist{language}
  \clearlist{translator}
  \clearfield{urldate}
  \clearfield{urlyear}
  \clearfield{urlmonth}
  \clearfield{pubstate}
  \clearfield{url}
}


%%%%%%%%%%%%%%%%%%%%%%%%%%%%%%%%%%%%%%%%%%%%%%%%%%%%%%%%%%%%%%%%%%%%%%%%%%%%%%
%%% 3. Cross-references
%%%%%%%%%%%%%%%%%%%%%%%%%%%%%%%%%%%%%%%%%%%%%%%%%%%%%%%%%%%%%%%%%%%%%%%%%%%%%%

% hyperref should be loaded late, but before zref-clever.
\usepackage[colorlinks=true,allcolors=blue]{hyperref}
\usepackage{zref-clever}
\usepackage{aliascnt}
\zcsetup{cap}

% cleveref-style spellings, so that sources written against \cref still work.
\providecommand{\cref}[1]{\zcref{#1}}
\providecommand{\Cref}[1]{\zcref{#1}}
\providecommand{\crefrange}[2]{\zcref{#1,#2}}
\providecommand{\Crefrange}[2]{\zcref{#1,#2}}
\providecommand{\labelcref}[1]{\zcref{#1}}

% Equations are referred to by number alone: (3.2), not "equation (3.2)".
\zcRefTypeSetup{equation}{
  name-sg=,
  Name-sg=,
  name-pl=,
  Name-pl=
}

% \NewTheoremWithZref{env}[shared counter]{Name}{Plural}
%
% With plain \newtheorem, every environment sharing the `theorem' counter gets
% referred to as "Theorem".  \aliascnt gives each one its own counter name so
% zref-clever can print "Lemma 3.2", "Definition 3.3", and so on.
\NewDocumentCommand{\NewTheoremWithZref}{ m o m m }{%
  \IfNoValueTF{#2}{%
    \newtheorem{#1}{#3}%
    \zcRefTypeSetup{#1}{%
      name-sg=#3, Name-sg=#3, name-pl=#4, Name-pl=#4%
    }%
  }{%
    \newaliascnt{#1}{#2}%
    \newtheorem{#1}[#1]{#3}%
    \aliascntresetthe{#1}%
    \zcRefTypeSetup{#1}{%
      name-sg=#3, Name-sg=#3, name-pl=#4, Name-pl=#4%
    }%
  }%
}


%%%%%%%%%%%%%%%%%%%%%%%%%%%%%%%%%%%%%%%%%%%%%%%%%%%%%%%%%%%%%%%%%%%%%%%%%%%%%%
%%% 4. Theorem environments
%%%%%%%%%%%%%%%%%%%%%%%%%%%%%%%%%%%%%%%%%%%%%%%%%%%%%%%%%%%%%%%%%%%%%%%%%%%%%%

% Results stated in the introduction: Theorem A, Corollary B, ...
\newtheorem{itheorem}{Theorem}
\renewcommand{\theitheorem}{\Alph{itheorem}}
\zcRefTypeSetup{itheorem}{
  name-sg=theorem, Name-sg=Theorem, name-pl=theorems, Name-pl=Theorems
}
\NewTheoremWithZref{icorollary}[itheorem]{Corollary}{Corollaries}
\NewTheoremWithZref{iproposition}[itheorem]{Proposition}{Propositions}

% Everything in the body shares one counter, numbered by section.
\newtheorem{theorem}{Theorem}[section]
\NewTheoremWithZref{proposition}[theorem]{Proposition}{Propositions}
\NewTheoremWithZref{corollary}[theorem]{Corollary}{Corollaries}
\NewTheoremWithZref{lemma}[theorem]{Lemma}{Lemmas}
\NewTheoremWithZref{assumption}[theorem]{Assumption}{Assumptions}
\NewTheoremWithZref{definition}[theorem]{Definition}{Definitions}
\NewTheoremWithZref{ques}[theorem]{Question}{Questions}
\NewTheoremWithZref{conjecture}[theorem]{Conjecture}{Conjectures}
\NewTheoremWithZref{algorithm}[theorem]{Algorithm}{Algorithms}

\theoremstyle{remark}
\NewTheoremWithZref{remark}[theorem]{Remark}{Remarks}
\NewTheoremWithZref{warning}[theorem]{Warning}{Warnings}

% Examples get their own per-section counter with roman numerals, and close
% with a club suit rather than a tombstone.
\NewTheoremWithZref{examplex}{Example}{Examples}
\makeatletter
\@addtoreset{examplex}{section}
\makeatother
\renewcommand{\theexamplex}{\arabic{section}.\roman{examplex}}
\newenvironment{example}
  {\pushQED{\qed}\renewcommand{\qedsymbol}{$\clubsuit$}\examplex}
  {\popQED\endexamplex}


%%%%%%%%%%%%%%%%%%%%%%%%%%%%%%%%%%%%%%%%%%%%%%%%%%%%%%%%%%%%%%%%%%%%%%%%%%%%%%
%%% 5. Equation numbering
%%%%%%%%%%%%%%%%%%%%%%%%%%%%%%%%%%%%%%%%%%%%%%%%%%%%%%%%%%%%%%%%%%%%%%%%%%%%%%

\numberwithin{equation}{section}
\renewcommand{\theequation}{\arabic{section}.\arabic{equation}}

% \newseq steps the equation counter once; each \subeqn inside the following
% *unnumbered* displays (\[...\], align*, gather*) then tags a line 3.4a, 3.4b,
% ....  Do not use \subeqn inside a numbered `equation', which would step the
% equation counter again and restart the letters.
\newcounter{subeqn}
\renewcommand{\thesubeqn}{\theequation\alph{subeqn}}
\makeatletter
% Freeze the number at the moment the counter steps.  amsart and friends are
% leqno classes, and there multline builds its tag box after amsmath has rolled
% back the counters from its measuring pass, so a bare \tag{\thesubeqn} prints
% one letter too low and leaves the hyperref destination out of step with it.
\newcommand{\subeqn}{\refstepcounter{subeqn}%
  \protected@edef\subeqn@tag{\thesubeqn}%
  \expandafter\tag\expandafter{\subeqn@tag}}
\@addtoreset{subeqn}{equation}
\makeatother
\newcommand{\newseq}{\refstepcounter{equation}}

\setcounter{tocdepth}{1}


%%%%%%%%%%%%%%%%%%%%%%%%%%%%%%%%%%%%%%%%%%%%%%%%%%%%%%%%%%%%%%%%%%%%%%%%%%%%%%
%%% 6. Index of notation
%%%%%%%%%%%%%%%%%%%%%%%%%%%%%%%%%%%%%%%%%%%%%%%%%%%%%%%%%%%%%%%%%%%%%%%%%%%%%%

% Set \anindexfalse to drop the index without touching the \notation calls
% scattered through the text (useful for journal submissions).
\newif\ifanindex
\anindextrue
\ifanindex
  \usepackage{anindex}                % local: anindex.sty
  \anindexsetup{
    heading,
    lines,
    title = {Index of Notation}
  }
\else
  % Same signature as anindex's \notation, so entries can stay in the source.
  \NewDocumentCommand{\notation}{ m m o }{}
  \NewDocumentCommand{\IndexOfNotation}{ }{}
\fi


%%%%%%%%%%%%%%%%%%%%%%%%%%%%%%%%%%%%%%%%%%%%%%%%%%%%%%%%%%%%%%%%%%%%%%%%%%%%%%
%%% 7. Standard macros
%%%%%%%%%%%%%%%%%%%%%%%%%%%%%%%%%%%%%%%%%%%%%%%%%%%%%%%%%%%%%%%%%%%%%%%%%%%%%%

\newcommand{\nc}{\newcommand}
\newcommand{\renc}{\renewcommand}

%%% Math alphabet families -------------------------------------------------
%
%   \cA ... \cZ   \mathcal        \fA ... \fZ, \fa ... \fz   \mathfrak
%   \sA ... \sZ   \mathscr        \BA ... \BZ, \Ba ... \Bz   \mathbf
%
% Generated in a loop rather than listed one per line.  Existing commands are
% never overwritten, so \renc is the way to repurpose one of these in a paper.
% Note \fi is a TeX primitive and is therefore skipped: write \mathfrak{i}.
\makeatletter
\newcommand{\bw@letterfamily}[3]{%
  \@tfor\bw@ltr:=#3\do{%
    \@ifundefined{#1\bw@ltr}{%
      \expandafter\edef\csname #1\bw@ltr\endcsname{\noexpand#2{\bw@ltr}}%
    }{}%
  }%
}
\bw@letterfamily{c}{\mathcal} {ABCDEFGHIJKLMNOPQRSTUVWXYZ}
\bw@letterfamily{s}{\mathscr} {ABCDEFGHIJKLMNOPQRSTUVWXYZ}
\bw@letterfamily{f}{\mathfrak}{ABCDEFGHIJKLMNOPQRSTUVWXYZ}
\bw@letterfamily{f}{\mathfrak}{abcdefghjklmnopqrstuvwxyz}
\bw@letterfamily{B}{\mathbf}  {ABCDEFGHIJKLMNOPQRSTUVWXYZ}
\bw@letterfamily{B}{\mathbf}  {abcdefghijklmnopqrstuvwxyz}
\makeatother

%%% Number systems ---------------------------------------------------------
\nc{\C}{\mathbb{C}}
\nc{\R}{\mathbb{R}}
\nc{\Z}{\mathbb{Z}}
\nc{\N}{\mathbb{N}}
\nc{\Q}{\mathbb{Q}}
\nc{\kk}{\Bbbk}                       % the base field

%%% Lie algebras -----------------------------------------------------------
\nc{\g}{\mathfrak{g}}
\nc{\h}{\mathfrak{h}}
\nc{\gl}{\mathfrak{gl}}
\nc{\fsl}{\mathfrak{sl}}
\nc{\fso}{\mathfrak{so}}
\nc{\fsp}{\mathfrak{sp}}
\nc{\pgl}{\mathfrak{pgl}}

%%% Greek shorthand --------------------------------------------------------
\nc{\al}{\alpha}      \nc{\be}{\beta}       \nc{\ga}{\gamma}
\nc{\de}{\delta}      \nc{\ep}{\epsilon}    \nc{\eps}{\varepsilon}
\nc{\la}{\lambda}     \nc{\si}{\sigma}      \nc{\om}{\omega}
\nc{\vp}{\varphi}     \nc{\La}{\Lambda}     \nc{\Th}{\Theta}

%%% Operators --------------------------------------------------------------
\DeclareMathOperator{\Hom}{Hom}
\DeclareMathOperator{\End}{End}
\DeclareMathOperator{\Ext}{Ext}
\DeclareMathOperator{\Tor}{Tor}
\DeclareMathOperator{\RHom}{RHom}
\DeclareMathOperator{\HOM}{HOM}
\DeclareMathOperator{\Aut}{Aut}
\DeclareMathOperator{\Fun}{Fun}
\DeclareMathOperator{\Sym}{Sym}
\DeclareMathOperator{\Spec}{Spec}
\DeclareMathOperator{\MaxSpec}{MaxSpec}
\DeclareMathOperator{\Proj}{Proj}
\DeclareMathOperator{\Frac}{Frac}
\DeclareMathOperator{\Coh}{Coh}
\DeclareMathOperator{\GL}{GL}
\DeclareMathOperator{\Mat}{Mat}
\DeclareMathOperator{\Kar}{Kar}
\DeclareMathOperator{\Add}{Add}
\DeclareMathOperator{\Ind}{Ind}
\DeclareMathOperator{\Res}{Res}
\DeclareMathOperator{\Tr}{Tr}
\DeclareMathOperator{\Id}{Id}
\DeclareMathOperator{\Ad}{Ad}
\DeclareMathOperator{\ad}{ad}
\DeclareMathOperator{\ch}{ch}
\DeclareMathOperator{\coker}{coker}
\DeclareMathOperator{\cosoc}{cosoc}
\DeclareMathOperator{\diag}{diag}
\DeclareMathOperator{\ev}{ev}
\DeclareMathOperator{\gr}{gr}
\DeclareMathOperator{\id}{id}
\DeclareMathOperator{\im}{im}
\DeclareMathOperator{\ind}{ind}
\DeclareMathOperator{\op}{op}
\DeclareMathOperator{\pr}{pr}
\DeclareMathOperator{\rad}{rad}
\DeclareMathOperator{\res}{res}
\DeclareMathOperator{\rev}{rev}
\DeclareMathOperator{\sgn}{sgn}
\DeclareMathOperator{\soc}{soc}
\DeclareMathOperator{\supp}{supp}
\DeclareMathOperator{\tr}{tr}
\DeclareMathOperator{\wt}{wt}
\nc{\mmod}{\operatorname{-mod}}       % as in $A\mmod$

%%% Assorted ---------------------------------------------------------------
\nc{\qbinom}[2]{\genfrac{[}{]}{0pt}{}{#1}{#2}}
\nc{\doubletilde}[1]{\tilde{\tilde{#1}}}
\nc{\arxiv}[1]{\href{http://arxiv.org/abs/#1}{\texttt{arXiv:\nolinkurl{#1}}}}


%%%%%%%%%%%%%%%%%%%%%%%%%%%%%%%%%%%%%%%%%%%%%%%%%%%%%%%%%%%%%%%%%%%%%%%%%%%%%%
%%% 8. Editorial macros
%%%%%%%%%%%%%%%%%%%%%%%%%%%%%%%%%%%%%%%%%%%%%%%%%%%%%%%%%%%%%%%%%%%%%%%%%%%%%%

% Margin/inline notes between coauthors.  Add one \nc per author, and switch
% \bwnotestrue to \bwnotesfalse to silence them all before circulating a draft.
\newif\ifbwnotes
\bwnotestrue
\ifbwnotes
  \nc{\Ben}[1]{{\color{red}[Ben: #1]}}
\else
  \nc{\Ben}[1]{}
\fi
 immediately after
%%% \documentclass.  Project-specific macros belong in main.tex (or their own
%%% file), not here: this file is meant to be copied unchanged between papers.
%%%
%%%   1. Packages
%%%   2. Bibliography (biblatex + biber)
%%%   3. Cross-references (zref-clever + aliascnt)
%%%   4. Theorem environments
%%%   5. Equation numbering
%%%   6. Index of notation (anindex)
%%%   7. Standard macros
%%%   8. Editorial macros


%%%%%%%%%%%%%%%%%%%%%%%%%%%%%%%%%%%%%%%%%%%%%%%%%%%%%%%%%%%%%%%%%%%%%%%%%%%%%%
%%% 1. Packages
%%%%%%%%%%%%%%%%%%%%%%%%%%%%%%%%%%%%%%%%%%%%%%%%%%%%%%%%%%%%%%%%%%%%%%%%%%%%%%

% These two must be set before bm loads; without them bm + stmaryrd together
% exhaust the 16 available math families.
\newcommand{\hmmax}{0}
\newcommand{\bmmax}{0}

\usepackage[english]{babel}
\usepackage{csquotes}
\usepackage[letterpaper,top=2cm,bottom=2cm,left=3cm,right=3cm,marginparwidth=1.75cm]{geometry}

\usepackage{amsmath,amsfonts,amssymb,mathtools}
\usepackage{mathrsfs}                 % \mathscr
\usepackage{bm,stmaryrd,nicefrac}
\usepackage{xcolor,graphicx,xspace}
\usepackage{xparse,enumitem}
\usepackage{tikz,tikz-cd}
\usepackage{quiver}                   % local: quiver.sty
\usepackage{tabularx}

% Uncomment while drafting to print label names in the margin.
% \usepackage[notref,notcite]{showkeys}


%%%%%%%%%%%%%%%%%%%%%%%%%%%%%%%%%%%%%%%%%%%%%%%%%%%%%%%%%%%%%%%%%%%%%%%%%%%%%%
%%% 2. Bibliography
%%%%%%%%%%%%%%%%%%%%%%%%%%%%%%%%%%%%%%%%%%%%%%%%%%%%%%%%%%%%%%%%%%%%%%%%%%%%%%

\usepackage[maxbibnames=99,backend=biber,url=false,style=alphabetic,doi=true,eprint=false]{biblatex}
\renewbibmacro*{url+urldate}{%
  % nothing here
}
\renewbibmacro{in:}{}
\AtEveryBibitem{
  \clearfield{isbn}
  \clearlist{location}
  \clearfield{month}
  \clearfield{day}
  \clearfield{series}
  \clearfield{pages}
  \clearlist{language}
  \clearlist{translator}
  \clearfield{urldate}
  \clearfield{urlyear}
  \clearfield{urlmonth}
  \clearfield{pubstate}
  \clearfield{url}
}


%%%%%%%%%%%%%%%%%%%%%%%%%%%%%%%%%%%%%%%%%%%%%%%%%%%%%%%%%%%%%%%%%%%%%%%%%%%%%%
%%% 3. Cross-references
%%%%%%%%%%%%%%%%%%%%%%%%%%%%%%%%%%%%%%%%%%%%%%%%%%%%%%%%%%%%%%%%%%%%%%%%%%%%%%

% hyperref should be loaded late, but before zref-clever.
\usepackage[colorlinks=true,allcolors=blue]{hyperref}
\usepackage{zref-clever}
\usepackage{aliascnt}
\zcsetup{cap}

% cleveref-style spellings, so that sources written against \cref still work.
\providecommand{\cref}[1]{\zcref{#1}}
\providecommand{\Cref}[1]{\zcref{#1}}
\providecommand{\crefrange}[2]{\zcref{#1,#2}}
\providecommand{\Crefrange}[2]{\zcref{#1,#2}}
\providecommand{\labelcref}[1]{\zcref{#1}}

% Equations are referred to by number alone: (3.2), not "equation (3.2)".
\zcRefTypeSetup{equation}{
  name-sg=,
  Name-sg=,
  name-pl=,
  Name-pl=
}

% \NewTheoremWithZref{env}[shared counter]{Name}{Plural}
%
% With plain \newtheorem, every environment sharing the `theorem' counter gets
% referred to as "Theorem".  \aliascnt gives each one its own counter name so
% zref-clever can print "Lemma 3.2", "Definition 3.3", and so on.
\NewDocumentCommand{\NewTheoremWithZref}{ m o m m }{%
  \IfNoValueTF{#2}{%
    \newtheorem{#1}{#3}%
    \zcRefTypeSetup{#1}{%
      name-sg=#3, Name-sg=#3, name-pl=#4, Name-pl=#4%
    }%
  }{%
    \newaliascnt{#1}{#2}%
    \newtheorem{#1}[#1]{#3}%
    \aliascntresetthe{#1}%
    \zcRefTypeSetup{#1}{%
      name-sg=#3, Name-sg=#3, name-pl=#4, Name-pl=#4%
    }%
  }%
}


%%%%%%%%%%%%%%%%%%%%%%%%%%%%%%%%%%%%%%%%%%%%%%%%%%%%%%%%%%%%%%%%%%%%%%%%%%%%%%
%%% 4. Theorem environments
%%%%%%%%%%%%%%%%%%%%%%%%%%%%%%%%%%%%%%%%%%%%%%%%%%%%%%%%%%%%%%%%%%%%%%%%%%%%%%

% Results stated in the introduction: Theorem A, Corollary B, ...
\newtheorem{itheorem}{Theorem}
\renewcommand{\theitheorem}{\Alph{itheorem}}
\zcRefTypeSetup{itheorem}{
  name-sg=theorem, Name-sg=Theorem, name-pl=theorems, Name-pl=Theorems
}
\NewTheoremWithZref{icorollary}[itheorem]{Corollary}{Corollaries}
\NewTheoremWithZref{iproposition}[itheorem]{Proposition}{Propositions}

% Everything in the body shares one counter, numbered by section.
\newtheorem{theorem}{Theorem}[section]
\NewTheoremWithZref{proposition}[theorem]{Proposition}{Propositions}
\NewTheoremWithZref{corollary}[theorem]{Corollary}{Corollaries}
\NewTheoremWithZref{lemma}[theorem]{Lemma}{Lemmas}
\NewTheoremWithZref{assumption}[theorem]{Assumption}{Assumptions}
\NewTheoremWithZref{definition}[theorem]{Definition}{Definitions}
\NewTheoremWithZref{ques}[theorem]{Question}{Questions}
\NewTheoremWithZref{conjecture}[theorem]{Conjecture}{Conjectures}
\NewTheoremWithZref{algorithm}[theorem]{Algorithm}{Algorithms}

\theoremstyle{remark}
\NewTheoremWithZref{remark}[theorem]{Remark}{Remarks}
\NewTheoremWithZref{warning}[theorem]{Warning}{Warnings}

% Examples get their own per-section counter with roman numerals, and close
% with a club suit rather than a tombstone.
\NewTheoremWithZref{examplex}{Example}{Examples}
\makeatletter
\@addtoreset{examplex}{section}
\makeatother
\renewcommand{\theexamplex}{\arabic{section}.\roman{examplex}}
\newenvironment{example}
  {\pushQED{\qed}\renewcommand{\qedsymbol}{$\clubsuit$}\examplex}
  {\popQED\endexamplex}


%%%%%%%%%%%%%%%%%%%%%%%%%%%%%%%%%%%%%%%%%%%%%%%%%%%%%%%%%%%%%%%%%%%%%%%%%%%%%%
%%% 5. Equation numbering
%%%%%%%%%%%%%%%%%%%%%%%%%%%%%%%%%%%%%%%%%%%%%%%%%%%%%%%%%%%%%%%%%%%%%%%%%%%%%%

\numberwithin{equation}{section}
\renewcommand{\theequation}{\arabic{section}.\arabic{equation}}

% \newseq steps the equation counter once; each \subeqn inside the following
% *unnumbered* displays (\[...\], align*, gather*) then tags a line 3.4a, 3.4b,
% ....  Do not use \subeqn inside a numbered `equation', which would step the
% equation counter again and restart the letters.
\newcounter{subeqn}
\renewcommand{\thesubeqn}{\theequation\alph{subeqn}}
\makeatletter
% Freeze the number at the moment the counter steps.  amsart and friends are
% leqno classes, and there multline builds its tag box after amsmath has rolled
% back the counters from its measuring pass, so a bare \tag{\thesubeqn} prints
% one letter too low and leaves the hyperref destination out of step with it.
\newcommand{\subeqn}{\refstepcounter{subeqn}%
  \protected@edef\subeqn@tag{\thesubeqn}%
  \expandafter\tag\expandafter{\subeqn@tag}}
\@addtoreset{subeqn}{equation}
\makeatother
\newcommand{\newseq}{\refstepcounter{equation}}

\setcounter{tocdepth}{1}


%%%%%%%%%%%%%%%%%%%%%%%%%%%%%%%%%%%%%%%%%%%%%%%%%%%%%%%%%%%%%%%%%%%%%%%%%%%%%%
%%% 6. Index of notation
%%%%%%%%%%%%%%%%%%%%%%%%%%%%%%%%%%%%%%%%%%%%%%%%%%%%%%%%%%%%%%%%%%%%%%%%%%%%%%

% Set \anindexfalse to drop the index without touching the \notation calls
% scattered through the text (useful for journal submissions).
\newif\ifanindex
\anindextrue
\ifanindex
  \usepackage{anindex}                % local: anindex.sty
  \anindexsetup{
    heading,
    lines,
    title = {Index of Notation}
  }
\else
  % Same signature as anindex's \notation, so entries can stay in the source.
  \NewDocumentCommand{\notation}{ m m o }{}
  \NewDocumentCommand{\IndexOfNotation}{ }{}
\fi


%%%%%%%%%%%%%%%%%%%%%%%%%%%%%%%%%%%%%%%%%%%%%%%%%%%%%%%%%%%%%%%%%%%%%%%%%%%%%%
%%% 7. Standard macros
%%%%%%%%%%%%%%%%%%%%%%%%%%%%%%%%%%%%%%%%%%%%%%%%%%%%%%%%%%%%%%%%%%%%%%%%%%%%%%

\newcommand{\nc}{\newcommand}
\newcommand{\renc}{\renewcommand}

%%% Math alphabet families -------------------------------------------------
%
%   \cA ... \cZ   \mathcal        \fA ... \fZ, \fa ... \fz   \mathfrak
%   \sA ... \sZ   \mathscr        \BA ... \BZ, \Ba ... \Bz   \mathbf
%
% Generated in a loop rather than listed one per line.  Existing commands are
% never overwritten, so \renc is the way to repurpose one of these in a paper.
% Note \fi is a TeX primitive and is therefore skipped: write \mathfrak{i}.
\makeatletter
\newcommand{\bw@letterfamily}[3]{%
  \@tfor\bw@ltr:=#3\do{%
    \@ifundefined{#1\bw@ltr}{%
      \expandafter\edef\csname #1\bw@ltr\endcsname{\noexpand#2{\bw@ltr}}%
    }{}%
  }%
}
\bw@letterfamily{c}{\mathcal} {ABCDEFGHIJKLMNOPQRSTUVWXYZ}
\bw@letterfamily{s}{\mathscr} {ABCDEFGHIJKLMNOPQRSTUVWXYZ}
\bw@letterfamily{f}{\mathfrak}{ABCDEFGHIJKLMNOPQRSTUVWXYZ}
\bw@letterfamily{f}{\mathfrak}{abcdefghjklmnopqrstuvwxyz}
\bw@letterfamily{B}{\mathbf}  {ABCDEFGHIJKLMNOPQRSTUVWXYZ}
\bw@letterfamily{B}{\mathbf}  {abcdefghijklmnopqrstuvwxyz}
\makeatother

%%% Number systems ---------------------------------------------------------
\nc{\C}{\mathbb{C}}
\nc{\R}{\mathbb{R}}
\nc{\Z}{\mathbb{Z}}
\nc{\N}{\mathbb{N}}
\nc{\Q}{\mathbb{Q}}
\nc{\kk}{\Bbbk}                       % the base field

%%% Lie algebras -----------------------------------------------------------
\nc{\g}{\mathfrak{g}}
\nc{\h}{\mathfrak{h}}
\nc{\gl}{\mathfrak{gl}}
\nc{\fsl}{\mathfrak{sl}}
\nc{\fso}{\mathfrak{so}}
\nc{\fsp}{\mathfrak{sp}}
\nc{\pgl}{\mathfrak{pgl}}

%%% Greek shorthand --------------------------------------------------------
\nc{\al}{\alpha}      \nc{\be}{\beta}       \nc{\ga}{\gamma}
\nc{\de}{\delta}      \nc{\ep}{\epsilon}    \nc{\eps}{\varepsilon}
\nc{\la}{\lambda}     \nc{\si}{\sigma}      \nc{\om}{\omega}
\nc{\vp}{\varphi}     \nc{\La}{\Lambda}     \nc{\Th}{\Theta}

%%% Operators --------------------------------------------------------------
\DeclareMathOperator{\Hom}{Hom}
\DeclareMathOperator{\End}{End}
\DeclareMathOperator{\Ext}{Ext}
\DeclareMathOperator{\Tor}{Tor}
\DeclareMathOperator{\RHom}{RHom}
\DeclareMathOperator{\HOM}{HOM}
\DeclareMathOperator{\Aut}{Aut}
\DeclareMathOperator{\Fun}{Fun}
\DeclareMathOperator{\Sym}{Sym}
\DeclareMathOperator{\Spec}{Spec}
\DeclareMathOperator{\MaxSpec}{MaxSpec}
\DeclareMathOperator{\Proj}{Proj}
\DeclareMathOperator{\Frac}{Frac}
\DeclareMathOperator{\Coh}{Coh}
\DeclareMathOperator{\GL}{GL}
\DeclareMathOperator{\Mat}{Mat}
\DeclareMathOperator{\Kar}{Kar}
\DeclareMathOperator{\Add}{Add}
\DeclareMathOperator{\Ind}{Ind}
\DeclareMathOperator{\Res}{Res}
\DeclareMathOperator{\Tr}{Tr}
\DeclareMathOperator{\Id}{Id}
\DeclareMathOperator{\Ad}{Ad}
\DeclareMathOperator{\ad}{ad}
\DeclareMathOperator{\ch}{ch}
\DeclareMathOperator{\coker}{coker}
\DeclareMathOperator{\cosoc}{cosoc}
\DeclareMathOperator{\diag}{diag}
\DeclareMathOperator{\ev}{ev}
\DeclareMathOperator{\gr}{gr}
\DeclareMathOperator{\id}{id}
\DeclareMathOperator{\im}{im}
\DeclareMathOperator{\ind}{ind}
\DeclareMathOperator{\op}{op}
\DeclareMathOperator{\pr}{pr}
\DeclareMathOperator{\rad}{rad}
\DeclareMathOperator{\res}{res}
\DeclareMathOperator{\rev}{rev}
\DeclareMathOperator{\sgn}{sgn}
\DeclareMathOperator{\soc}{soc}
\DeclareMathOperator{\supp}{supp}
\DeclareMathOperator{\tr}{tr}
\DeclareMathOperator{\wt}{wt}
\nc{\mmod}{\operatorname{-mod}}       % as in $A\mmod$

%%% Assorted ---------------------------------------------------------------
\nc{\qbinom}[2]{\genfrac{[}{]}{0pt}{}{#1}{#2}}
\nc{\doubletilde}[1]{\tilde{\tilde{#1}}}
\nc{\arxiv}[1]{\href{http://arxiv.org/abs/#1}{\texttt{arXiv:\nolinkurl{#1}}}}


%%%%%%%%%%%%%%%%%%%%%%%%%%%%%%%%%%%%%%%%%%%%%%%%%%%%%%%%%%%%%%%%%%%%%%%%%%%%%%
%%% 8. Editorial macros
%%%%%%%%%%%%%%%%%%%%%%%%%%%%%%%%%%%%%%%%%%%%%%%%%%%%%%%%%%%%%%%%%%%%%%%%%%%%%%

% Margin/inline notes between coauthors.  Add one \nc per author, and switch
% \bwnotestrue to \bwnotesfalse to silence them all before circulating a draft.
\newif\ifbwnotes
\bwnotestrue
\ifbwnotes
  \nc{\Ben}[1]{{\color{red}[Ben: #1]}}
\else
  \nc{\Ben}[1]{}
\fi
 immediately after
%%% \documentclass.  Project-specific macros belong in main.tex (or their own
%%% file), not here: this file is meant to be copied unchanged between papers.
%%%
%%%   1. Packages
%%%   2. Bibliography (biblatex + biber)
%%%   3. Cross-references (zref-clever + aliascnt)
%%%   4. Theorem environments
%%%   5. Equation numbering
%%%   6. Index of notation (anindex)
%%%   7. Standard macros
%%%   8. Editorial macros


%%%%%%%%%%%%%%%%%%%%%%%%%%%%%%%%%%%%%%%%%%%%%%%%%%%%%%%%%%%%%%%%%%%%%%%%%%%%%%
%%% 1. Packages
%%%%%%%%%%%%%%%%%%%%%%%%%%%%%%%%%%%%%%%%%%%%%%%%%%%%%%%%%%%%%%%%%%%%%%%%%%%%%%

% These two must be set before bm loads; without them bm + stmaryrd together
% exhaust the 16 available math families.
\newcommand{\hmmax}{0}
\newcommand{\bmmax}{0}

\usepackage[english]{babel}
\usepackage{csquotes}
\usepackage[letterpaper,top=2cm,bottom=2cm,left=3cm,right=3cm,marginparwidth=1.75cm]{geometry}

\usepackage{amsmath,amsfonts,amssymb,mathtools}
\usepackage{mathrsfs}                 % \mathscr
\usepackage{bm,stmaryrd,nicefrac}
\usepackage{xcolor,graphicx,xspace}
\usepackage{xparse,enumitem}
\usepackage{tikz,tikz-cd}
\usepackage{quiver}                   % local: quiver.sty
\usepackage{tabularx}

% Uncomment while drafting to print label names in the margin.
% \usepackage[notref,notcite]{showkeys}


%%%%%%%%%%%%%%%%%%%%%%%%%%%%%%%%%%%%%%%%%%%%%%%%%%%%%%%%%%%%%%%%%%%%%%%%%%%%%%
%%% 2. Bibliography
%%%%%%%%%%%%%%%%%%%%%%%%%%%%%%%%%%%%%%%%%%%%%%%%%%%%%%%%%%%%%%%%%%%%%%%%%%%%%%

\usepackage[maxbibnames=99,backend=biber,url=false,style=alphabetic,doi=true,eprint=false]{biblatex}
\renewbibmacro*{url+urldate}{%
  % nothing here
}
\renewbibmacro{in:}{}
\AtEveryBibitem{
  \clearfield{isbn}
  \clearlist{location}
  \clearfield{month}
  \clearfield{day}
  \clearfield{series}
  \clearfield{pages}
  \clearlist{language}
  \clearlist{translator}
  \clearfield{urldate}
  \clearfield{urlyear}
  \clearfield{urlmonth}
  \clearfield{pubstate}
  \clearfield{url}
}


%%%%%%%%%%%%%%%%%%%%%%%%%%%%%%%%%%%%%%%%%%%%%%%%%%%%%%%%%%%%%%%%%%%%%%%%%%%%%%
%%% 3. Cross-references
%%%%%%%%%%%%%%%%%%%%%%%%%%%%%%%%%%%%%%%%%%%%%%%%%%%%%%%%%%%%%%%%%%%%%%%%%%%%%%

% hyperref should be loaded late, but before zref-clever.
\usepackage[colorlinks=true,allcolors=blue]{hyperref}
\usepackage{zref-clever}
\usepackage{aliascnt}
\zcsetup{cap}

% cleveref-style spellings, so that sources written against \cref still work.
\providecommand{\cref}[1]{\zcref{#1}}
\providecommand{\Cref}[1]{\zcref{#1}}
\providecommand{\crefrange}[2]{\zcref{#1,#2}}
\providecommand{\Crefrange}[2]{\zcref{#1,#2}}
\providecommand{\labelcref}[1]{\zcref{#1}}

% Equations are referred to by number alone: (3.2), not "equation (3.2)".
\zcRefTypeSetup{equation}{
  name-sg=,
  Name-sg=,
  name-pl=,
  Name-pl=
}

% \NewTheoremWithZref{env}[shared counter]{Name}{Plural}
%
% With plain \newtheorem, every environment sharing the `theorem' counter gets
% referred to as "Theorem".  \aliascnt gives each one its own counter name so
% zref-clever can print "Lemma 3.2", "Definition 3.3", and so on.
\NewDocumentCommand{\NewTheoremWithZref}{ m o m m }{%
  \IfNoValueTF{#2}{%
    \newtheorem{#1}{#3}%
    \zcRefTypeSetup{#1}{%
      name-sg=#3, Name-sg=#3, name-pl=#4, Name-pl=#4%
    }%
  }{%
    \newaliascnt{#1}{#2}%
    \newtheorem{#1}[#1]{#3}%
    \aliascntresetthe{#1}%
    \zcRefTypeSetup{#1}{%
      name-sg=#3, Name-sg=#3, name-pl=#4, Name-pl=#4%
    }%
  }%
}


%%%%%%%%%%%%%%%%%%%%%%%%%%%%%%%%%%%%%%%%%%%%%%%%%%%%%%%%%%%%%%%%%%%%%%%%%%%%%%
%%% 4. Theorem environments
%%%%%%%%%%%%%%%%%%%%%%%%%%%%%%%%%%%%%%%%%%%%%%%%%%%%%%%%%%%%%%%%%%%%%%%%%%%%%%

% Results stated in the introduction: Theorem A, Corollary B, ...
\newtheorem{itheorem}{Theorem}
\renewcommand{\theitheorem}{\Alph{itheorem}}
\zcRefTypeSetup{itheorem}{
  name-sg=theorem, Name-sg=Theorem, name-pl=theorems, Name-pl=Theorems
}
\NewTheoremWithZref{icorollary}[itheorem]{Corollary}{Corollaries}
\NewTheoremWithZref{iproposition}[itheorem]{Proposition}{Propositions}

% Everything in the body shares one counter, numbered by section.
\newtheorem{theorem}{Theorem}[section]
\NewTheoremWithZref{proposition}[theorem]{Proposition}{Propositions}
\NewTheoremWithZref{corollary}[theorem]{Corollary}{Corollaries}
\NewTheoremWithZref{lemma}[theorem]{Lemma}{Lemmas}
\NewTheoremWithZref{assumption}[theorem]{Assumption}{Assumptions}
\NewTheoremWithZref{definition}[theorem]{Definition}{Definitions}
\NewTheoremWithZref{ques}[theorem]{Question}{Questions}
\NewTheoremWithZref{conjecture}[theorem]{Conjecture}{Conjectures}
\NewTheoremWithZref{algorithm}[theorem]{Algorithm}{Algorithms}

\theoremstyle{remark}
\NewTheoremWithZref{remark}[theorem]{Remark}{Remarks}
\NewTheoremWithZref{warning}[theorem]{Warning}{Warnings}

% Examples get their own per-section counter with roman numerals, and close
% with a club suit rather than a tombstone.
\NewTheoremWithZref{examplex}{Example}{Examples}
\makeatletter
\@addtoreset{examplex}{section}
\makeatother
\renewcommand{\theexamplex}{\arabic{section}.\roman{examplex}}
\newenvironment{example}
  {\pushQED{\qed}\renewcommand{\qedsymbol}{$\clubsuit$}\examplex}
  {\popQED\endexamplex}


%%%%%%%%%%%%%%%%%%%%%%%%%%%%%%%%%%%%%%%%%%%%%%%%%%%%%%%%%%%%%%%%%%%%%%%%%%%%%%
%%% 5. Equation numbering
%%%%%%%%%%%%%%%%%%%%%%%%%%%%%%%%%%%%%%%%%%%%%%%%%%%%%%%%%%%%%%%%%%%%%%%%%%%%%%

\numberwithin{equation}{section}
\renewcommand{\theequation}{\arabic{section}.\arabic{equation}}

% \newseq steps the equation counter once; each \subeqn inside the following
% *unnumbered* displays (\[...\], align*, gather*) then tags a line 3.4a, 3.4b,
% ....  Do not use \subeqn inside a numbered `equation', which would step the
% equation counter again and restart the letters.
\newcounter{subeqn}
\renewcommand{\thesubeqn}{\theequation\alph{subeqn}}
\makeatletter
% Freeze the number at the moment the counter steps.  amsart and friends are
% leqno classes, and there multline builds its tag box after amsmath has rolled
% back the counters from its measuring pass, so a bare \tag{\thesubeqn} prints
% one letter too low and leaves the hyperref destination out of step with it.
\newcommand{\subeqn}{\refstepcounter{subeqn}%
  \protected@edef\subeqn@tag{\thesubeqn}%
  \expandafter\tag\expandafter{\subeqn@tag}}
\@addtoreset{subeqn}{equation}
\makeatother
\newcommand{\newseq}{\refstepcounter{equation}}

\setcounter{tocdepth}{1}


%%%%%%%%%%%%%%%%%%%%%%%%%%%%%%%%%%%%%%%%%%%%%%%%%%%%%%%%%%%%%%%%%%%%%%%%%%%%%%
%%% 6. Index of notation
%%%%%%%%%%%%%%%%%%%%%%%%%%%%%%%%%%%%%%%%%%%%%%%%%%%%%%%%%%%%%%%%%%%%%%%%%%%%%%

% Set \anindexfalse to drop the index without touching the \notation calls
% scattered through the text (useful for journal submissions).
\newif\ifanindex
\anindextrue
\ifanindex
  \usepackage{anindex}                % local: anindex.sty
  \anindexsetup{
    heading,
    lines,
    title = {Index of Notation}
  }
\else
  % Same signature as anindex's \notation, so entries can stay in the source.
  \NewDocumentCommand{\notation}{ m m o }{}
  \NewDocumentCommand{\IndexOfNotation}{ }{}
\fi


%%%%%%%%%%%%%%%%%%%%%%%%%%%%%%%%%%%%%%%%%%%%%%%%%%%%%%%%%%%%%%%%%%%%%%%%%%%%%%
%%% 7. Standard macros
%%%%%%%%%%%%%%%%%%%%%%%%%%%%%%%%%%%%%%%%%%%%%%%%%%%%%%%%%%%%%%%%%%%%%%%%%%%%%%

\newcommand{\nc}{\newcommand}
\newcommand{\renc}{\renewcommand}

%%% Math alphabet families -------------------------------------------------
%
%   \cA ... \cZ   \mathcal        \fA ... \fZ, \fa ... \fz   \mathfrak
%   \sA ... \sZ   \mathscr        \BA ... \BZ, \Ba ... \Bz   \mathbf
%
% Generated in a loop rather than listed one per line.  Existing commands are
% never overwritten, so \renc is the way to repurpose one of these in a paper.
% Note \fi is a TeX primitive and is therefore skipped: write \mathfrak{i}.
\makeatletter
\newcommand{\bw@letterfamily}[3]{%
  \@tfor\bw@ltr:=#3\do{%
    \@ifundefined{#1\bw@ltr}{%
      \expandafter\edef\csname #1\bw@ltr\endcsname{\noexpand#2{\bw@ltr}}%
    }{}%
  }%
}
\bw@letterfamily{c}{\mathcal} {ABCDEFGHIJKLMNOPQRSTUVWXYZ}
\bw@letterfamily{s}{\mathscr} {ABCDEFGHIJKLMNOPQRSTUVWXYZ}
\bw@letterfamily{f}{\mathfrak}{ABCDEFGHIJKLMNOPQRSTUVWXYZ}
\bw@letterfamily{f}{\mathfrak}{abcdefghjklmnopqrstuvwxyz}
\bw@letterfamily{B}{\mathbf}  {ABCDEFGHIJKLMNOPQRSTUVWXYZ}
\bw@letterfamily{B}{\mathbf}  {abcdefghijklmnopqrstuvwxyz}
\makeatother

%%% Number systems ---------------------------------------------------------
\nc{\C}{\mathbb{C}}
\nc{\R}{\mathbb{R}}
\nc{\Z}{\mathbb{Z}}
\nc{\N}{\mathbb{N}}
\nc{\Q}{\mathbb{Q}}
\nc{\kk}{\Bbbk}                       % the base field

%%% Lie algebras -----------------------------------------------------------
\nc{\g}{\mathfrak{g}}
\nc{\h}{\mathfrak{h}}
\nc{\gl}{\mathfrak{gl}}
\nc{\fsl}{\mathfrak{sl}}
\nc{\fso}{\mathfrak{so}}
\nc{\fsp}{\mathfrak{sp}}
\nc{\pgl}{\mathfrak{pgl}}

%%% Greek shorthand --------------------------------------------------------
\nc{\al}{\alpha}      \nc{\be}{\beta}       \nc{\ga}{\gamma}
\nc{\de}{\delta}      \nc{\ep}{\epsilon}    \nc{\eps}{\varepsilon}
\nc{\la}{\lambda}     \nc{\si}{\sigma}      \nc{\om}{\omega}
\nc{\vp}{\varphi}     \nc{\La}{\Lambda}     \nc{\Th}{\Theta}

%%% Operators --------------------------------------------------------------
\DeclareMathOperator{\Hom}{Hom}
\DeclareMathOperator{\End}{End}
\DeclareMathOperator{\Ext}{Ext}
\DeclareMathOperator{\Tor}{Tor}
\DeclareMathOperator{\RHom}{RHom}
\DeclareMathOperator{\HOM}{HOM}
\DeclareMathOperator{\Aut}{Aut}
\DeclareMathOperator{\Fun}{Fun}
\DeclareMathOperator{\Sym}{Sym}
\DeclareMathOperator{\Spec}{Spec}
\DeclareMathOperator{\MaxSpec}{MaxSpec}
\DeclareMathOperator{\Proj}{Proj}
\DeclareMathOperator{\Frac}{Frac}
\DeclareMathOperator{\Coh}{Coh}
\DeclareMathOperator{\GL}{GL}
\DeclareMathOperator{\Mat}{Mat}
\DeclareMathOperator{\Kar}{Kar}
\DeclareMathOperator{\Add}{Add}
\DeclareMathOperator{\Ind}{Ind}
\DeclareMathOperator{\Res}{Res}
\DeclareMathOperator{\Tr}{Tr}
\DeclareMathOperator{\Id}{Id}
\DeclareMathOperator{\Ad}{Ad}
\DeclareMathOperator{\ad}{ad}
\DeclareMathOperator{\ch}{ch}
\DeclareMathOperator{\coker}{coker}
\DeclareMathOperator{\cosoc}{cosoc}
\DeclareMathOperator{\diag}{diag}
\DeclareMathOperator{\ev}{ev}
\DeclareMathOperator{\gr}{gr}
\DeclareMathOperator{\id}{id}
\DeclareMathOperator{\im}{im}
\DeclareMathOperator{\ind}{ind}
\DeclareMathOperator{\op}{op}
\DeclareMathOperator{\pr}{pr}
\DeclareMathOperator{\rad}{rad}
\DeclareMathOperator{\res}{res}
\DeclareMathOperator{\rev}{rev}
\DeclareMathOperator{\sgn}{sgn}
\DeclareMathOperator{\soc}{soc}
\DeclareMathOperator{\supp}{supp}
\DeclareMathOperator{\tr}{tr}
\DeclareMathOperator{\wt}{wt}
\nc{\mmod}{\operatorname{-mod}}       % as in $A\mmod$

%%% Assorted ---------------------------------------------------------------
\nc{\qbinom}[2]{\genfrac{[}{]}{0pt}{}{#1}{#2}}
\nc{\doubletilde}[1]{\tilde{\tilde{#1}}}
\nc{\arxiv}[1]{\href{http://arxiv.org/abs/#1}{\texttt{arXiv:\nolinkurl{#1}}}}


%%%%%%%%%%%%%%%%%%%%%%%%%%%%%%%%%%%%%%%%%%%%%%%%%%%%%%%%%%%%%%%%%%%%%%%%%%%%%%
%%% 8. Editorial macros
%%%%%%%%%%%%%%%%%%%%%%%%%%%%%%%%%%%%%%%%%%%%%%%%%%%%%%%%%%%%%%%%%%%%%%%%%%%%%%

% Margin/inline notes between coauthors.  Add one \nc per author, and switch
% \bwnotestrue to \bwnotesfalse to silence them all before circulating a draft.
\newif\ifbwnotes
\bwnotestrue
\ifbwnotes
  \nc{\Ben}[1]{{\color{red}[Ben: #1]}}
\else
  \nc{\Ben}[1]{}
\fi
 immediately after
%%% \documentclass.  Project-specific macros belong in main.tex (or their own
%%% file), not here: this file is meant to be copied unchanged between papers.
%%%
%%%   1. Packages
%%%   2. Bibliography (biblatex + biber)
%%%   3. Cross-references (zref-clever + aliascnt)
%%%   4. Theorem environments
%%%   5. Equation numbering
%%%   6. Index of notation (anindex)
%%%   7. Standard macros
%%%   8. Editorial macros


%%%%%%%%%%%%%%%%%%%%%%%%%%%%%%%%%%%%%%%%%%%%%%%%%%%%%%%%%%%%%%%%%%%%%%%%%%%%%%
%%% 1. Packages
%%%%%%%%%%%%%%%%%%%%%%%%%%%%%%%%%%%%%%%%%%%%%%%%%%%%%%%%%%%%%%%%%%%%%%%%%%%%%%

% These two must be set before bm loads; without them bm + stmaryrd together
% exhaust the 16 available math families.
\newcommand{\hmmax}{0}
\newcommand{\bmmax}{0}

\usepackage[english]{babel}
\usepackage{csquotes}
\usepackage[letterpaper,top=2cm,bottom=2cm,left=3cm,right=3cm,marginparwidth=1.75cm]{geometry}

\usepackage{amsmath,amsfonts,amssymb,mathtools}
\usepackage{mathrsfs}                 % \mathscr
\usepackage{bm,stmaryrd,nicefrac}
\usepackage{xcolor,graphicx,xspace}
\usepackage{xparse,enumitem}
\usepackage{tikz,tikz-cd}
\usepackage{quiver}                   % local: quiver.sty
\usepackage{tabularx}

% Uncomment while drafting to print label names in the margin.
% \usepackage[notref,notcite]{showkeys}


%%%%%%%%%%%%%%%%%%%%%%%%%%%%%%%%%%%%%%%%%%%%%%%%%%%%%%%%%%%%%%%%%%%%%%%%%%%%%%
%%% 2. Bibliography
%%%%%%%%%%%%%%%%%%%%%%%%%%%%%%%%%%%%%%%%%%%%%%%%%%%%%%%%%%%%%%%%%%%%%%%%%%%%%%

\usepackage[maxbibnames=99,backend=biber,url=false,style=alphabetic,doi=true,eprint=false]{biblatex}
\renewbibmacro*{url+urldate}{%
  % nothing here
}
\renewbibmacro{in:}{}
\AtEveryBibitem{
  \clearfield{isbn}
  \clearlist{location}
  \clearfield{month}
  \clearfield{day}
  \clearfield{series}
  \clearfield{pages}
  \clearlist{language}
  \clearlist{translator}
  \clearfield{urldate}
  \clearfield{urlyear}
  \clearfield{urlmonth}
  \clearfield{pubstate}
  \clearfield{url}
}


%%%%%%%%%%%%%%%%%%%%%%%%%%%%%%%%%%%%%%%%%%%%%%%%%%%%%%%%%%%%%%%%%%%%%%%%%%%%%%
%%% 3. Cross-references
%%%%%%%%%%%%%%%%%%%%%%%%%%%%%%%%%%%%%%%%%%%%%%%%%%%%%%%%%%%%%%%%%%%%%%%%%%%%%%

% hyperref should be loaded late, but before zref-clever.
\usepackage[colorlinks=true,allcolors=blue]{hyperref}
\usepackage{zref-clever}
\usepackage{aliascnt}
\zcsetup{cap}

% cleveref-style spellings, so that sources written against \cref still work.
\providecommand{\cref}[1]{\zcref{#1}}
\providecommand{\Cref}[1]{\zcref{#1}}
\providecommand{\crefrange}[2]{\zcref{#1,#2}}
\providecommand{\Crefrange}[2]{\zcref{#1,#2}}
\providecommand{\labelcref}[1]{\zcref{#1}}

% Equations are referred to by number alone: (3.2), not "equation (3.2)".
\zcRefTypeSetup{equation}{
  name-sg=,
  Name-sg=,
  name-pl=,
  Name-pl=
}

% \NewTheoremWithZref{env}[shared counter]{Name}{Plural}
%
% With plain \newtheorem, every environment sharing the `theorem' counter gets
% referred to as "Theorem".  \aliascnt gives each one its own counter name so
% zref-clever can print "Lemma 3.2", "Definition 3.3", and so on.
\NewDocumentCommand{\NewTheoremWithZref}{ m o m m }{%
  \IfNoValueTF{#2}{%
    \newtheorem{#1}{#3}%
    \zcRefTypeSetup{#1}{%
      name-sg=#3, Name-sg=#3, name-pl=#4, Name-pl=#4%
    }%
  }{%
    \newaliascnt{#1}{#2}%
    \newtheorem{#1}[#1]{#3}%
    \aliascntresetthe{#1}%
    \zcRefTypeSetup{#1}{%
      name-sg=#3, Name-sg=#3, name-pl=#4, Name-pl=#4%
    }%
  }%
}


%%%%%%%%%%%%%%%%%%%%%%%%%%%%%%%%%%%%%%%%%%%%%%%%%%%%%%%%%%%%%%%%%%%%%%%%%%%%%%
%%% 4. Theorem environments
%%%%%%%%%%%%%%%%%%%%%%%%%%%%%%%%%%%%%%%%%%%%%%%%%%%%%%%%%%%%%%%%%%%%%%%%%%%%%%

% Results stated in the introduction: Theorem A, Corollary B, ...
\newtheorem{itheorem}{Theorem}
\renewcommand{\theitheorem}{\Alph{itheorem}}
\zcRefTypeSetup{itheorem}{
  name-sg=theorem, Name-sg=Theorem, name-pl=theorems, Name-pl=Theorems
}
\NewTheoremWithZref{icorollary}[itheorem]{Corollary}{Corollaries}
\NewTheoremWithZref{iproposition}[itheorem]{Proposition}{Propositions}

% Everything in the body shares one counter, numbered by section.
\newtheorem{theorem}{Theorem}[section]
\NewTheoremWithZref{proposition}[theorem]{Proposition}{Propositions}
\NewTheoremWithZref{corollary}[theorem]{Corollary}{Corollaries}
\NewTheoremWithZref{lemma}[theorem]{Lemma}{Lemmas}
\NewTheoremWithZref{assumption}[theorem]{Assumption}{Assumptions}
\NewTheoremWithZref{definition}[theorem]{Definition}{Definitions}
\NewTheoremWithZref{ques}[theorem]{Question}{Questions}
\NewTheoremWithZref{conjecture}[theorem]{Conjecture}{Conjectures}
\NewTheoremWithZref{algorithm}[theorem]{Algorithm}{Algorithms}

\theoremstyle{remark}
\NewTheoremWithZref{remark}[theorem]{Remark}{Remarks}
\NewTheoremWithZref{warning}[theorem]{Warning}{Warnings}

% Examples get their own per-section counter with roman numerals, and close
% with a club suit rather than a tombstone.
\NewTheoremWithZref{examplex}{Example}{Examples}
\makeatletter
\@addtoreset{examplex}{section}
\makeatother
\renewcommand{\theexamplex}{\arabic{section}.\roman{examplex}}
\newenvironment{example}
  {\pushQED{\qed}\renewcommand{\qedsymbol}{$\clubsuit$}\examplex}
  {\popQED\endexamplex}


%%%%%%%%%%%%%%%%%%%%%%%%%%%%%%%%%%%%%%%%%%%%%%%%%%%%%%%%%%%%%%%%%%%%%%%%%%%%%%
%%% 5. Equation numbering
%%%%%%%%%%%%%%%%%%%%%%%%%%%%%%%%%%%%%%%%%%%%%%%%%%%%%%%%%%%%%%%%%%%%%%%%%%%%%%

\numberwithin{equation}{section}
\renewcommand{\theequation}{\arabic{section}.\arabic{equation}}

% \newseq steps the equation counter once; each \subeqn inside the following
% *unnumbered* displays (\[...\], align*, gather*) then tags a line 3.4a, 3.4b,
% ....  Do not use \subeqn inside a numbered `equation', which would step the
% equation counter again and restart the letters.
\newcounter{subeqn}
\renewcommand{\thesubeqn}{\theequation\alph{subeqn}}
\makeatletter
% Freeze the number at the moment the counter steps.  amsart and friends are
% leqno classes, and there multline builds its tag box after amsmath has rolled
% back the counters from its measuring pass, so a bare \tag{\thesubeqn} prints
% one letter too low and leaves the hyperref destination out of step with it.
\newcommand{\subeqn}{\refstepcounter{subeqn}%
  \protected@edef\subeqn@tag{\thesubeqn}%
  \expandafter\tag\expandafter{\subeqn@tag}}
\@addtoreset{subeqn}{equation}
\makeatother
\newcommand{\newseq}{\refstepcounter{equation}}

\setcounter{tocdepth}{1}


%%%%%%%%%%%%%%%%%%%%%%%%%%%%%%%%%%%%%%%%%%%%%%%%%%%%%%%%%%%%%%%%%%%%%%%%%%%%%%
%%% 6. Index of notation
%%%%%%%%%%%%%%%%%%%%%%%%%%%%%%%%%%%%%%%%%%%%%%%%%%%%%%%%%%%%%%%%%%%%%%%%%%%%%%

% Set \anindexfalse to drop the index without touching the \notation calls
% scattered through the text (useful for journal submissions).
\newif\ifanindex
\anindextrue
\ifanindex
  \usepackage{anindex}                % local: anindex.sty
  \anindexsetup{
    heading,
    lines,
    title = {Index of Notation}
  }
\else
  % Same signature as anindex's \notation, so entries can stay in the source.
  \NewDocumentCommand{\notation}{ m m o }{}
  \NewDocumentCommand{\IndexOfNotation}{ }{}
\fi


%%%%%%%%%%%%%%%%%%%%%%%%%%%%%%%%%%%%%%%%%%%%%%%%%%%%%%%%%%%%%%%%%%%%%%%%%%%%%%
%%% 7. Standard macros
%%%%%%%%%%%%%%%%%%%%%%%%%%%%%%%%%%%%%%%%%%%%%%%%%%%%%%%%%%%%%%%%%%%%%%%%%%%%%%

\newcommand{\nc}{\newcommand}
\newcommand{\renc}{\renewcommand}

%%% Math alphabet families -------------------------------------------------
%
%   \cA ... \cZ   \mathcal        \fA ... \fZ, \fa ... \fz   \mathfrak
%   \sA ... \sZ   \mathscr        \BA ... \BZ, \Ba ... \Bz   \mathbf
%
% Generated in a loop rather than listed one per line.  Existing commands are
% never overwritten, so \renc is the way to repurpose one of these in a paper.
% Note \fi is a TeX primitive and is therefore skipped: write \mathfrak{i}.
\makeatletter
\newcommand{\bw@letterfamily}[3]{%
  \@tfor\bw@ltr:=#3\do{%
    \@ifundefined{#1\bw@ltr}{%
      \expandafter\edef\csname #1\bw@ltr\endcsname{\noexpand#2{\bw@ltr}}%
    }{}%
  }%
}
\bw@letterfamily{c}{\mathcal} {ABCDEFGHIJKLMNOPQRSTUVWXYZ}
\bw@letterfamily{s}{\mathscr} {ABCDEFGHIJKLMNOPQRSTUVWXYZ}
\bw@letterfamily{f}{\mathfrak}{ABCDEFGHIJKLMNOPQRSTUVWXYZ}
\bw@letterfamily{f}{\mathfrak}{abcdefghjklmnopqrstuvwxyz}
\bw@letterfamily{B}{\mathbf}  {ABCDEFGHIJKLMNOPQRSTUVWXYZ}
\bw@letterfamily{B}{\mathbf}  {abcdefghijklmnopqrstuvwxyz}
\makeatother

%%% Number systems ---------------------------------------------------------
\nc{\C}{\mathbb{C}}
\nc{\R}{\mathbb{R}}
\nc{\Z}{\mathbb{Z}}
\nc{\N}{\mathbb{N}}
\nc{\Q}{\mathbb{Q}}
\nc{\kk}{\Bbbk}                       % the base field

%%% Lie algebras -----------------------------------------------------------
\nc{\g}{\mathfrak{g}}
\nc{\h}{\mathfrak{h}}
\nc{\gl}{\mathfrak{gl}}
\nc{\fsl}{\mathfrak{sl}}
\nc{\fso}{\mathfrak{so}}
\nc{\fsp}{\mathfrak{sp}}
\nc{\pgl}{\mathfrak{pgl}}

%%% Greek shorthand --------------------------------------------------------
\nc{\al}{\alpha}      \nc{\be}{\beta}       \nc{\ga}{\gamma}
\nc{\de}{\delta}      \nc{\ep}{\epsilon}    \nc{\eps}{\varepsilon}
\nc{\la}{\lambda}     \nc{\si}{\sigma}      \nc{\om}{\omega}
\nc{\vp}{\varphi}     \nc{\La}{\Lambda}     \nc{\Th}{\Theta}

%%% Operators --------------------------------------------------------------
\DeclareMathOperator{\Hom}{Hom}
\DeclareMathOperator{\End}{End}
\DeclareMathOperator{\Ext}{Ext}
\DeclareMathOperator{\Tor}{Tor}
\DeclareMathOperator{\RHom}{RHom}
\DeclareMathOperator{\HOM}{HOM}
\DeclareMathOperator{\Aut}{Aut}
\DeclareMathOperator{\Fun}{Fun}
\DeclareMathOperator{\Sym}{Sym}
\DeclareMathOperator{\Spec}{Spec}
\DeclareMathOperator{\MaxSpec}{MaxSpec}
\DeclareMathOperator{\Proj}{Proj}
\DeclareMathOperator{\Frac}{Frac}
\DeclareMathOperator{\Coh}{Coh}
\DeclareMathOperator{\GL}{GL}
\DeclareMathOperator{\Mat}{Mat}
\DeclareMathOperator{\Kar}{Kar}
\DeclareMathOperator{\Add}{Add}
\DeclareMathOperator{\Ind}{Ind}
\DeclareMathOperator{\Res}{Res}
\DeclareMathOperator{\Tr}{Tr}
\DeclareMathOperator{\Id}{Id}
\DeclareMathOperator{\Ad}{Ad}
\DeclareMathOperator{\ad}{ad}
\DeclareMathOperator{\ch}{ch}
\DeclareMathOperator{\coker}{coker}
\DeclareMathOperator{\cosoc}{cosoc}
\DeclareMathOperator{\diag}{diag}
\DeclareMathOperator{\ev}{ev}
\DeclareMathOperator{\gr}{gr}
\DeclareMathOperator{\id}{id}
\DeclareMathOperator{\im}{im}
\DeclareMathOperator{\ind}{ind}
\DeclareMathOperator{\op}{op}
\DeclareMathOperator{\pr}{pr}
\DeclareMathOperator{\rad}{rad}
\DeclareMathOperator{\res}{res}
\DeclareMathOperator{\rev}{rev}
\DeclareMathOperator{\sgn}{sgn}
\DeclareMathOperator{\soc}{soc}
\DeclareMathOperator{\supp}{supp}
\DeclareMathOperator{\tr}{tr}
\DeclareMathOperator{\wt}{wt}
\nc{\mmod}{\operatorname{-mod}}       % as in $A\mmod$

%%% Assorted ---------------------------------------------------------------
\nc{\qbinom}[2]{\genfrac{[}{]}{0pt}{}{#1}{#2}}
\nc{\doubletilde}[1]{\tilde{\tilde{#1}}}
\nc{\arxiv}[1]{\href{http://arxiv.org/abs/#1}{\texttt{arXiv:\nolinkurl{#1}}}}


%%%%%%%%%%%%%%%%%%%%%%%%%%%%%%%%%%%%%%%%%%%%%%%%%%%%%%%%%%%%%%%%%%%%%%%%%%%%%%
%%% 8. Editorial macros
%%%%%%%%%%%%%%%%%%%%%%%%%%%%%%%%%%%%%%%%%%%%%%%%%%%%%%%%%%%%%%%%%%%%%%%%%%%%%%

% Margin/inline notes between coauthors.  Add one \nc per author, and switch
% \bwnotestrue to \bwnotesfalse to silence them all before circulating a draft.
\newif\ifbwnotes
\bwnotestrue
\ifbwnotes
  \nc{\Ben}[1]{{\color{red}[Ben: #1]}}
\else
  \nc{\Ben}[1]{}
\fi
